\chapter{Duration, the arrow of time, and the boundary between potential and actual}\label{ch:time}

\section{Three kinds of time}

\emph{The Two Faces of Time}~\cite{MahaffeyTime} separates three notions that are often run together.
\begin{description}
\item[Coordinate time] is a label on the manifold. In the relational tradition of Barbour
      and Bertotti it is redundant, and the dynamics can be written without
      it~\cite{BarbourBertotti1982,Barbour1994}.
\item[Proper time] $\tau=\int\sqrt{-g_{\mu\nu}dx^\mu dx^\nu}/c$ is path length along a
      worldline. It is geometric and, in principle, reversible. On a null geodesic
      $d\tau=0$ exactly.
\item[Duration] is the accumulated count of irreversible actualization events. It is
      thermodynamic, it has a direction, and in IAM it is measured by $E(a)$.
\end{description}
Photons accumulate neither proper time nor duration. Massive particles accumulate both, but
only duration is irreversible. The arrow of time is then not a statistical statement about
low-entropy initial conditions~\cite{Penrose1989,Carroll2010}. It is the direction in which
the information ledger of the horizon is written, and every entry costs $\kB T_H\ln2$.
\conjecture

\section{Before and after the electroweak transition}

Before electroweak symmetry breaking ($t\approx9\times10^{-12}$~s; Chapter~\ref{ch:electroweak}) every Standard Model particle is
massless. There are no timelike worldlines carrying rest mass, no gravitational decoherence
of massive bodies, and $E(a)\equiv0$. After the transition the matter sector exists and the
photon remains massless under exact $U(1)_{\mathrm{em}}$, so the two sectors of
Chapter~\ref{ch:dual} separate permanently. The source papers leave open whether the
selection of a vacuum at the transition was itself the first actualization event, and this
edition leaves it open too.

\begin{remark}
The activation function of Eq.~\eqref{part2:eq:Ea} is normalized to $E(1)=1$ and is exponentially
small, not zero, at $a_{\mathrm{EW}}\sim10^{-15}$: $E=\exp(1-10^{15})$. ``$E\equiv0$ before
the transition'' is therefore a physical statement about the sector content, not a property
of the formula, which has no feature at $a_{\mathrm{EW}}$.
\end{remark}

\section{The boundary between potential and actual}

\emph{The Boundary Between Potential and Actual}~\cite{MahaffeyBoundary} uses the virial
partition of Chapter~\ref{ch:virial_law}. When a self-gravitating structure virializes, half of
the released binding energy, the ``geometric half'', is deposited in curvature. The other
half, the ``informational half'', is the Landauer cost of actualization and must be written
to an encoding surface. On galactic scales and above, the paper argues, the only available
surface is a black-hole horizon. Where no such horizon can form, the partition cannot close
and the structure disperses.

\prediction\ This gives a minimum velocity dispersion, $\sigma_{\mathrm{crit}}\approx4\ \mathrm{km\,s^{-1}}$, below which dark-matter halos do
not virialize, and hence a floor on the internal dispersions of satellite galaxies. The derivation of the 4~km\,s$^{-1}$ figure is not given in the
source papers and is not reproduced here. \openprob

The data already bear on it. Segue~2 is identified as a galaxy, not a star cluster, by the spread of its stars' metallicities ($-2.85$ to
$-1.33$ in [Fe/H]), and its velocity dispersion is unresolved: $\sigma_v<2.2$~km\,s$^{-1}$ at 90\,\% confidence and $<2.6$ at 95\,\%, from 25
member stars~\cite{Kirby2013Segue2}. \observed{} Other ultra-faint dwarfs reach 2--4~km\,s$^{-1}$~\cite{Simon2011Segue,Simon2019}. Across the Milky Way's
satellites, 25 of 54 have measured dispersions below 4~km\,s$^{-1}$ today, Segue~2 among them. \observed{} The floor as stated is rejected. The
only form left to test is a floor on the dispersion at infall, before tides stripped the satellites, which needs each satellite's orbital history.
\openprob

\section{Dark matter as the geometric half}

The same papers identify dark matter with the accumulated geometric half of past
virialization events: already actual, and not writing new information. This identification
is what justifies the baryon fraction $\Omega_b/\Omega_m$ in Chapter~\ref{ch:lambda}.
\conjecture\ It is a large claim, and it has to face every phenomenological success of
particle dark matter at once: the CMB peak heights, the Bullet Cluster offset, and
structure formation before recombination. It is recorded here as the framework's position,
not as a result.
